% Discussion — Paper 5 (base-drag closure intercomparison).
% Body only; master supplies \section{Discussion}.
% Every number per paper/_shared/CANONICAL_FACTS.md (sections A, B, C, D).
% Headline: base-drag fidelity is the highest-leverage modeling investment for
% integrated supersonic-rocket apogee; the shock pre-pass is essential for
% LOCAL-flow fidelity but apogee-inert (0.15 pp); the in-sample FINNED_BASE_K is
% handled by the holdout and bounded by external benchmarks. Relate to (do NOT
% repeat) the companion JSR integrated result. Discuss generality across vehicle
% classes. No Devan-Ashwood; Cd_base = 0.064 + 0.186/M^2 is empirical, anchored on
% NACA TN 3393 + ESDU 77021. \cite{yu2026jsr} is the companion-paper placeholder key.

\subsection{Base-drag fidelity is the highest-leverage modeling investment}
\label{sec:disc-budget}

The central result of the ablation (Section~\ref{sec:results}) is a sharp
ordering of mechanisms by their leverage on integrated apogee. Disabling the
finned-base augmentation closure moves the corpus mean apogee error by $8.10$~pp,
an order of magnitude more than any other supersonic mechanism: Van~Driest~II
compressible skin friction contributes $0.87$~pp~\cite{hopkins1971}, the
DATCOM~4.1.5.1 fin wave-drag term $0.39$~pp~\cite{datcom1978,ulmann1956}, and the
shock-geometry pre-pass only $0.15$~pp. The Pitts--Nielsen--Kaattari fin-body
interference floor moves apogee by $0.00$~pp. The base-flow closure is therefore
not merely one term among several: on this corpus it is the single dominant
driver of where the integrated trajectory ends up.

This ordering carries a direct prescription for where modeling effort should go.
A practitioner extending an open-source trajectory code to supersonic speeds
faces a long menu of candidate refinements --- higher-order wave-drag theory,
boundary-layer transition modeling, crossflow corrections, refined pitch
damping --- and finite time to implement them. The ablation says plainly that,
for the quantity most users ultimately care about, base-flow fidelity is the
highest-leverage investment, returning roughly nine times the apogee accuracy per
unit of modeling effort as the next-ranked mechanism and about 54 times that of
the shock-geometry pre-pass. The closure branches described here --- the
Chapman--Korst turbulent shear-layer theory~\cite{chapman1955,chapman1958}, the
Chapman laminar branch~\cite{chapman1950}, the ESDU~77021 bodies-of-revolution
correlation~\cite{esdu77021}, the empirical $C_{d,\text{base}} = 0.064 +
0.186/M^2$ form, and the Viswanath boattail-relief
correction~\cite{viswanath1996} --- together make up the term that carries this
leverage. We do not isolate the apogee contribution of each branch separately;
the ablation attributes the leverage to the base-drag term as a whole, and within
it the finned-base augmentation is the component the corpus is most sensitive to.
The prescription is that refining base-flow fidelity, and especially the
treatment of the finned base, returns more apogee accuracy than refining any
other supersonic sub-model first.

\subsection{Mechanism attribution versus local-flow fidelity}
\label{sec:disc-prepass}

It is tempting to equate architectural prominence with apogee leverage, and the
shock-geometry pre-pass is the case where that equation breaks. The pre-pass is
the most structurally celebrated component of the supersonic extension: it solves
the oblique-shock and Taylor--Maccoll cone field once per timestep and threads the
post-shock local conditions into every downstream model. Yet it moves integrated
apogee by only $0.15$~pp, with a maximum single-flight effect of $3.6$~pp
(FMJ Black~Rock~6, $M=2.46$). The reason is not that the pre-pass is wrong but
that apogee integrates a trajectory whose drag impulse is dominated by the
lower-Mach portion of the ascent, and the pre-pass is inert below $M=1$. A
mechanism can be both physically essential and apogee-irrelevant.

The pre-pass earns its place for a different reason: local-flow fidelity. The
post-shock conditions it produces feed fin loads and static stability, where they
have been verified to match Taylor--Maccoll cone
tables~\cite{taylormaccoll1933,sims1964} to better than $1\%$ and the analytic
shoulder-expansion result to numerical precision. It is also the architectural
seam that makes the downstream supersonic closures --- including the base-flow
models studied here, which depend on the local edge conditions at separation ---
well-posed in the first place. The honest framing is thus twofold: the pre-pass
is indispensable for correct \emph{local} aerodynamics (fin loads, static and
dynamic stability, recovery-event timing) and as an enabling architecture, but it
is not a gross-apogee mechanism. Conflating the two would misdirect a reader
toward refining shock geometry when the apogee error budget sits almost entirely
in the base-flow term.

This is precisely why mechanism attribution, rather than aggregate accuracy
alone, is the useful product of an ablation study. An integrated parity statistic
tells a user that the tool works; it does not tell a developer \emph{which}
sub-model to improve next. The decomposition here answers that second question,
and its answer --- base flow first, everything else second --- redirects future
modeling and validation effort toward the term that actually controls the
trajectory-integrated result.

\subsection{In-sample exposure of the finned-base closure, stated plainly}
\label{sec:disc-insample}

The dominant mechanism is also the one carrying empirical calibration, and
intellectual honesty requires confronting that directly rather than burying it.
There are two distinct in-sample exposures here, applied by different parts of the
code, and they must be separated to be answered honestly.

The first exposure is in the base-flow \emph{scale} constants. Two corpus-frozen
boundary-layer scale factors enter the base-drag treatment ---
$\mathrm{THICK\_BL\_K}=2.2$, anchored on the Raven flight, and
$\mathrm{SLENDER\_BODY\_K}=0.0025$, anchored on Raven, Rabia, and Kinsel. Because
those two constants were fitted with reference to flights that also appear in the
validation corpus, the headline aggregate is \emph{partly in-sample}. This is the
exposure the decontaminated holdout was designed to neutralize, and it neutralizes
it directly. Every flight that either scale constant touched --- Raven, Rabia,
Rabia Short~Fin~Can, Kinsel, and the Torrent diagnostic case --- was placed in the
development partition, leaving a genuinely blind holdout. The development partition
($n=13$) returns a mean signed error of $+0.22\%$ and a mean absolute error of
$5.47\%$; the holdout ($n=12$), which neither scale constant ever saw, returns mean
$-1.03\%$ and MAE $3.95\%$. The holdout is \emph{more} accurate than the
development set. Overfit closures degrade on unseen data; a closure whose scale
factors improve on the holdout is generalizing the underlying base-flow physics
rather than memorizing the calibration flights. This inversion --- holdout better
than dev --- is strong evidence that the two scale constants generalize.

The second exposure is the finned-base augmentation factor itself
($\mathrm{FINNED\_BASE\_K}$), the constant whose disabling produces the $8.10$~pp
dominant-mechanism result. We state plainly that the decontaminated holdout above
does \emph{not} isolate this factor: the holdout split was constructed to quarantine
the two boundary-layer scale constants, not the flights that informed the
finned-base augmentation, so it cannot be read as a clean-room test of the dominant
lever. The defense of $\mathrm{FINNED\_BASE\_K}$ therefore rests not on that holdout
but on external bounding. The component benchmarks of Section~\ref{sec:results}
fix the relevant physics from outside the corpus entirely: the turbulent clean-body
closures reproduce the NACA~TN~3393 nonlifting-body base pressures to $15.9\%$ MAPE
(and the Hart free-flight subset to $4.0\%$)~\cite{chapman1955}, the Chapman laminar
branch to $4.4\%$~\cite{chapman1955,chapman1950}, and the magnitude of the
finned-base augmentation itself sits inside the $40$--$60\%$ over-clean-body band
reported in aeroballistic-range testing of the Basic Finner~\cite{dupuis1997} and in
Hoerner's afterbody-drag compendium~\cite{hoerner1965}. The augmentation factor is
thus fixed by external finned-afterbody data, not by the flight corpus; it does not
invent the mechanism. We make no claim that any of these constants are physically
derived --- they are disclosed calibrations --- but the external Basic-Finner anchor
of the augmentation factor, together with the holdout generalization of the two
corpus-frozen scale constants, is what licenses naming the finned base the dominant
apogee driver. Elevating $\mathrm{FINNED\_BASE\_K}$ from its present total-drag
anchor to an isolated finned-base base-pressure benchmark is a clean direction for
future work.

\subsection{Relation to the integrated companion study}
\label{sec:disc-companion}

The present paper and its integrated-validation companion answer different
questions, and the divergence is the point. The question here is one of mechanism
leverage --- \emph{which} closure governs the trajectory-integrated result --- and
its answer is that the apogee accuracy rests disproportionately on the base-flow
closure while the celebrated shock-geometry architecture is, by the apogee metric,
nearly inert. The companion paper~\cite{yu2026jsr} asks instead whether the
integrated tool predicts supersonic apogee and how it compares to the commercial
baseline~\cite{rogers2015}, and reports that integrated-parity result; it does not
decompose \emph{why} the prediction lands where it does. The two studies share a
corpus and a code base but pursue distinct ends, and the value of reading them
together is that the present decomposition explains the structure underlying the
companion's aggregate. The claim defended here is a statement about mechanism
leverage, not about aggregate parity, and it stands on the ablation and the
holdout alone.

\subsection{Generality across vehicle classes}
\label{sec:disc-generality}

How far does ``base drag governs apogee'' carry beyond the present corpus? The
mechanism that makes base drag dominant --- a large, separated aft-face pressure
deficit acting through the low-supersonic band where the climb-to-apogee integral
is weighted most heavily --- is generic to slender, fin-stabilized,
flat-or-near-flat-base vehicles, the class that spans the corpus from minimum-diameter
amateur high-power rockets through university-competition airframes to the
single- and two-stage sounding rockets that reach the upper Mach range. Within
that class the conclusion is expected to be robust, because every member shares
the same qualitative base-flow topology and the same trajectory weighting toward
the lower-supersonic regime. The corpus itself spans Mach~$0.54$ to $4.33$ and a
wide spread of fineness ratios and fin geometries, and the dominance of the
drag-coefficient term is not an artifact of one airframe: a per-flight
sensitivity sweep (Fig.~\ref{fig:tornado}) shows that the overall drag-coefficient
scale --- of which base drag is the dominant supersonic component --- is the
single most leveraged parameter for predicted apogee across structurally diverse
vehicles spanning the low-supersonic to upper-Mach range, dwarfing launch
geometry and integration-timestep sensitivities on every flight shown.

\begin{figure}[htbp]
\centering
\includegraphics[width=\linewidth]{sensitivity_tornado_composite.png}
\caption{Per-flight apogee sensitivity (tornado) for representative
high-supersonic vehicles: change in predicted apogee for a $\pm 1\sigma$ swing in
each model parameter. The overall drag-coefficient scale --- whose supersonic
budget is dominated by the base-drag term identified in
Table~\ref{tab:ablation} --- is the most leveraged parameter on every vehicle,
far ahead of launch-rod angle, launch altitude, and integration timestep. This
flight-by-flight consistency across structurally different airframes is direct
evidence that the leverage of the base-drag term on integrated apogee is generic
to the slender fin-stabilized class, not an artifact of a single geometry.}
\label{fig:tornado}
\end{figure}

The conclusion does, however, have boundaries that should be stated rather than
assumed. Three vehicle attributes shift the apogee error budget away from base
drag. First, \emph{aggressive boattailing}: a well-designed aft taper converts a
large fraction of the base area to a recompressing surface, so the Viswanath
relief~\cite{viswanath1996} shrinks the base term and a larger share of the
budget migrates to forebody and fin wave drag --- and, as the retained RM-10
negative benchmark shows, an upstream boattail can drive the finned-base
augmentation outside its flat-base calibration envelope entirely. Second,
\emph{boundary-layer state}: the corpus is overwhelmingly turbulent at
separation, but a true perfect-finish, laminar-to-the-base vehicle is governed by
the Chapman laminar branch~\cite{chapman1950}, whose much smaller base drag (a
$44\%$ component swing relative to the turbulent form on the same TN~3393 data)
reduces the term's absolute apogee leverage even though base flow remains the
relevant closure. Third, \emph{flight regime}: the dominance is a low-supersonic
phenomenon. For a vehicle that spends most of its impulse subsonically, base drag
recedes behind skin friction and form drag; for a sustained-hypersonic vehicle,
forebody wave drag and real-gas effects grow in relative weight as the
$0.186/M^2$ term decays, and the present corpus offers no inferential support
above Mach~$4.33$. The defensible scope of the result is therefore: \emph{for
slender, fin-stabilized, blunt-or-lightly-tapered-base vehicles flying through the
transonic and low-supersonic regime, the base-drag closure --- and specifically
the treatment of the finned base --- is the highest-leverage determinant of
integrated apogee.} That is the regime in which the overwhelming majority of
amateur, collegiate, and meteorological sounding-rocket flights operate, which is
what makes the result broadly useful rather than corpus-specific; it is also a
falsifiable boundary, and extending the ablation to boattailed, perfect-finish,
and sustained-hypersonic vehicles is a natural test of where it ceases to hold.
